\documentclass[a4paper,12pt]{article}
\usepackage[T1]{fontenc}
\usepackage[latin9]{inputenc}

\setlength{\parindent}{3mm}
\usepackage{setspace}
\onehalfspacing

\usepackage{amsmath}
\usepackage{url}
\usepackage{verbatim}
\usepackage[numbers]{natbib}
\usepackage{nomencl}
% the following is useful when we have the old nomencl.sty package
\providecommand{\printnomenclature}{\printglossary}
\providecommand{\makenomenclature}{\makeglossary}
\usepackage{varioref}
\usepackage{prettyref}
\usepackage{subscript}

\usepackage{graphicx}

%%%%%%%%%%%%%%%%%%%%%%%%%%%%%% LyX specific LaTeX commands.
\newcommand{\lyxline}[1][1pt]{%
  \par\noindent%
  \rule[.5ex]{\linewidth}{#1}\par}

\newcommand{\lyxarrow}{\leavevmode\,$\triangleright$\,\allowbreak}

\begin{document}

\tableofcontents

\noindent This paragraph is not indented.

\section{References\index{References}}

Let's start with simple things: a label~\label{lab:test} and a
reference~\ref{lab:test}; note that I have tested ``unbreakable
space'' without warning (and worse than that, I just tested english
quotes too...).

Of course there are other kind of references, like page
reference~\pageref{lab:test}, but also equation
reference~\eqref{lab:test} (from amsmath package), or varioref's
equivalents~\vref{lab:test} and~\vpageref{lab:test}, without
forgetting pretty references like~\prettyref{lab:test}.

\section{Cites\index{Cites}}

Let's start with simple things: a \textbackslash{}cite: \cite[after]{article-crossref}
and two \textbackslash{}cites: \cite[after]{whole-set,article-crossref}

Now the natbib things:

\textbackslash{}citet: \citet[before][after]{article-crossref} \textbackslash{}citet{*}:
\citet*[before][after]{article-crossref} \textbackslash{}Citet: \Citet[before][after]{article-crossref}
\textbackslash{}Citet{*}: \Citet*[before][after]{article-crossref}

\textbackslash{}citep: \citep[before][after]{article-crossref} \textbackslash{}citep{*}:
\citep*[before][after]{article-crossref} \textbackslash{}Citep: \Citep[before][after]{article-crossref}
\textbackslash{}Citep{*}: \Citep*[before][after]{article-crossref}

\textbackslash{}citealt: \citealt[before][after]{article-crossref}
\textbackslash{}citealt{*}: \citealt*[before][after]{article-crossref}
\textbackslash{}Citealt: \Citealt[before][after]{article-crossref}
\textbackslash{}Citealt{*}: \Citealt*[before][after]{article-crossref}

\textbackslash{}citealp: \citealp[before][after]{article-crossref}
\textbackslash{}citealp{*}: \citealp*[before][after]{article-crossref}
\textbackslash{}Citealp: \Citep[before][after]{article-crossref}
\textbackslash{}Citealp{*}: \Citealp*[before][after]{article-crossref}

\textbackslash{}citeauthor: \citeauthor[before][after]{article-crossref}
\textbackslash{}citeauthor{*}: \citeauthor*[before][after]{article-crossref}
\textbackslash{}Citeauthor: \Citeauthor[before][after]{article-crossref}
\textbackslash{}Citeauthor{*}: \Citeauthor*[before][after]{article-crossref}

\textbackslash{}citeyear: \citeyear[before][after]{article-crossref}

\textbackslash{}citeyearpar: \citeyearpar[before][after]{article-crossref} 

\textbackslash{}nocite: \nocite{article-crossref}

\bibliographystyle{unsrt}
\bibliography{xampl}

\section{Input files\index{Input files}}

We can input files too, like this 
\subsection{\label{sub:External-Subsection}External Subsection}

This is a small dummy child document to show how files can be inserted
to another document. Here are some accented characters to make sure
the encoding is passed to included files: ��
, or with the include
variant 
\subsection{\label{sub:External-Subsection}External Subsection}

This is a small dummy child document to show how files can be inserted
to another document. Here are some accented characters to make sure
the encoding is passed to included files: ��


If you prefer verbatim input, you can choose
between~\verbatiminput{foo} or~\verbatiminput*{foo}.

\section{URL\index{URL}\nomenclature[www]{URL}{uniform resource locator}}

An URL: \url{http://www.lyx.org}
\nomenclature{URL2}{uniform resource locator}

\section{Lists\index{Lists}}

\listoffigures
\listoftables

\section{Graphics\index{Graphics}}

There is also some basic support for graphics, in the form
\includegraphics{foo.eps}, or the slightly more elaborate
\includegraphics[bb=10bp 0bp 96bp 96bp,clip,height=1cm, width=1cm]{foo.eps}.

\section{Special formattings\index{Special formattings}}

\subsection{LyX line}

test

{\tiny \lyxline{\tiny}}{\tiny \par}

test {\Huge Test} {\tiny test} test

\lyxline{\normalsize}

test

{\Huge \lyxline{\Huge}}{\Huge \par}

test

\subsection{Line breaks}

They can also or be broken by a newline\\
or a newline command \newline
or by a line break \linebreak
or by a defined line break \linebreak[4]

There are even newlines with weird arguments, but these are not
handled by LyX\\*[1cm]
so we try to use ERT in this case.

\subsection{Page breaks}

They can also or be broken by a newpage \newpage
or by a page break \pagebreak
or by a defined page break \pagebreak[4]

\section{Special characters\index{Special characters}}

Then one has those macros with a long name for a short meaning, like
\textasciitilde, \textasciicircum{} or \textbackslash{}, and the characters
that LaTeX wants to espace because they are active, like \_\&\#\$\{\}\%.

And what about special characters like hyphe\-nation mark,
ellipsis\ldots, and end-of-sentence\@. LyX also supports a menu
separator\lyxarrow{}and a spif\textcompwordmark{}fy ligature break.

A sub\textsubscript{sc\emph{ript}} and super\textsuperscript{script
with $a^2+b^2=c^2$ math}.

\printindex{}

\printnomenclature{}

\end{document}
